\section{Domain Analysis}
\subsection{Application Context and Utilized Technologies}
The primary objective of this laboratory work is to implement a \textbf{preemptive multitasking
system} using \textbf{FreeRTOS} on the ESP32 microcontroller. This approach extends the bare-metal
cooperative scheduler from Laboratory Work 3 by introducing a real-time operating system that
provides true preemptive scheduling, inter-task synchronization primitives, and drift-corrected
periodic execution.

The application uses FreeRTOS \textbf{preemptive tasks} managed by the RTOS kernel scheduler.
Each task runs as an independent execution context with its own stack. Periodic execution is
achieved through \texttt{vTaskDelayUntil}, which provides precise, drift-corrected intervals.
An initial \texttt{vTaskDelay} is used to apply the startup offset before the periodic loop begins,
mirroring the offset field of the bare-metal \texttt{TaskContext} structure.

Task synchronization is achieved through two FreeRTOS primitives:
\begin{itemize}
  \item A \textbf{binary semaphore} used by Task 1 to signal Task 2 when a new press event occurs
  \item A \textbf{mutex} used by Tasks 1, 2, and 3 to protect shared statistical variables from
    concurrent access
\end{itemize}

Communication with the user is achieved through the standard \texttt{printf()} function, which on
ESP32 is natively routed to the UART interface without requiring the \texttt{fdevopen} redirection
needed on AVR.

\subsection{Hardware and Software Components}

\begin{center}
\begin{tabular}{| p{4cm} | p{10cm} |}
\hline
\textbf{Component} & \textbf{Description \& Role} \\ \hline
\textbf{ESP32 DevKit} & Dual-core microcontroller board with Xtensa LX6 that runs FreeRTOS and
  all application tasks preemptively. \\ \hline
\textbf{Push Button} & Input peripheral connected to GPIO\,15 with internal pull-up resistor.
  Used to detect press/release transitions for duration measurement. \\ \hline
\textbf{Green LED} & Connected to GPIO\,13. Lights up to indicate a short button press ($< 500$ ms). \\ \hline
\textbf{Red LED} & Connected to GPIO\,26. Lights up to indicate a long button press ($\geq 500$ ms). \\ \hline
\textbf{Yellow LED} & Connected to GPIO\,32. Performs rapid blinking to acknowledge each press
  (5 blinks for short, 10 for long). \\ \hline
\textbf{$220\,\Omega$ Resistors} & Current limiting resistors for each LED to keep current within
  safe operating limits. \\ \hline
\textbf{Serial-to-USB Interface} & Facilitates UART communication between the MCU and the terminal
  for periodic report output at 115200 baud. \\ \hline
\textbf{FreeRTOS} & Real-time operating system built into the ESP32 Arduino core. Provides
  preemptive task scheduling, binary semaphores, and mutexes. \\ \hline
\textbf{Neovim + Arduino-nvim plugin} & Development environment with \texttt{arduino-cli} integration
  for compilation, upload, and monitoring. \\ \hline
\textbf{Breadboard} & Prototyping base for interconnecting the button, LEDs, resistors, and the
  ESP32 without soldering. \\ \hline
\end{tabular}
\end{center}

\subsection{System Architecture and Solution Justification}
The system adopts a \textbf{preemptive FreeRTOS scheduler architecture}. Unlike the cooperative
bare-metal design from Laboratory Work 3, FreeRTOS can interrupt any running task when a higher-priority
task becomes ready, making the system more responsive and realistic for concurrent embedded applications.

The key architectural insight is that the \textbf{task functions themselves remain unchanged} -- they
are still plain, one-shot functions with no internal loops. The periodicity and offset logic is
encapsulated entirely inside the scheduler service, which wraps each task function in a FreeRTOS
\texttt{periodicTaskRunner}:

\begin{lstlisting}
static void periodicTaskRunner(void *pvParameters) {
  TaskWrapper *ctx = (TaskWrapper *)pvParameters;
  vTaskDelay(pdMS_TO_TICKS(ctx->offsetMs));        // apply offset once
  TickType_t xLastWakeTime = xTaskGetTickCount();
  for (;;) {
    ctx->func();                                     // call plain task function
    vTaskDelayUntil(&xLastWakeTime, pdMS_TO_TICKS(ctx->periodMs));
  }
}
\end{lstlisting}

This design minimizes changes between the bare-metal and FreeRTOS versions, satisfying the
portability requirement of the laboratory work.

The data flow is structured as follows:
\begin{enumerate}
  \item \textbf{FreeRTOS Scheduler:} Preemptively runs Task 1 every 10\,ms, Task 2 every 50\,ms,
    and Task 3 every 10\,seconds based on their configured periods and priorities.
  \item \textbf{Task 1:} Detects button transitions, writes shared press data under mutex, lights
    the appropriate LED, and gives the binary semaphore.
  \item \textbf{Task 2:} Non-blockingly checks the semaphore. If taken, reads press data under mutex,
    updates statistics under mutex, and drives the yellow LED blink sequence.
  \item \textbf{Task 3:} Acquires the mutex to atomically snapshot and reset all counters, then prints
    the report via \texttt{printf()}.
\end{enumerate}

\subsection{Relevant Case Study}
A real-world application of preemptive multitasking with synchronization is found in
\textbf{industrial sensor monitoring systems} with safety requirements. Unlike cooperative systems,
preemptive scheduling ensures that high-priority tasks (e.g., emergency stop detection) can
immediately interrupt lower-priority tasks. FreeRTOS mutexes prevent data corruption when multiple
tasks share sensor readings or counters -- a critical requirement in IEC\,61508-compliant safety
systems. The binary semaphore pattern used here (producer signals consumer without busy-waiting)
is a standard technique in event-driven embedded firmware.
